This section records the design decisions and invariants of the CadmIA
implementation of IA-DEVS that apply to every experiment in this repository.

\subsection*{Foundational References}

The experiments in this repository reproduce models from the literature;
each experiment section cites its source.  UA-DEVS and IA-DEVS are
defined in~\cite{VDW21}, which is the primary reference for all formalism
definitions used here.

% ─────────────────────────────────────────────────────────────────────────────
\subsection*{Interval Representation}

IA-DEVS replaces every scalar quantity with an \emph{interval}.  CadmIA
represents an interval as a pair $[\ell, u]$ where $\ell$ and $u$ are values
of the same underlying scalar type.  Any interval kind is permitted:

\begin{center}
\begin{tabular}{@{}ll@{}}
  \toprule
  Kind & Condition \\
  \midrule
  Closed         & $\ell \leq u$                         \\
  Open           & $\ell < u$ (endpoints excluded)       \\
  Half-open      & one endpoint excluded                 \\
  Point          & $\ell = u$ (point interval)           \\
  Empty          & $\ell > u$ (no element)               \\
  \bottomrule
\end{tabular}
\end{center}

There is no requirement that intervals be closed.  The choice of interval kind
is part of the model specification for each experiment and should be justified
there.

% ─────────────────────────────────────────────────────────────────────────────
\subsection*{TIME Type}

The interval type for time is \texttt{interval<T>} for any scalar \texttt{T}.
Any type that is a valid TIME type in Cadmium or CDBoost is also a valid
\texttt{T}: fixed-point rational types (\texttt{decimal<N>}), integer types, and
any user-defined totally-ordered arithmetic type.

Floating-point types (\texttt{float}, \texttt{double}) are known to introduce
causality errors in Cadmium and CDBoost when period values are not exactly
representable.  In CadmIA, \texttt{float} and \texttt{double} are also valid
\texttt{T} choices, and \emph{they produce correct IA-DEVS results} provided the
arithmetic obeys the outward-rounding convention: lower-bound computations use
round-toward-$-\infty$ and upper-bound computations use round-toward-$+\infty$.
This is not the IEEE~754 default mode; see the Floating-Point paragraph of the
Outward-Rounding Invariant section for the required setup.

The choice of \texttt{T} is an approximation decision recorded in each
experiment's spec.

% ─────────────────────────────────────────────────────────────────────────────
\subsection*{Outward-Rounding Invariant}

The fundamental soundness requirement for IA-DEVS is that every interval
must \emph{contain} the exact value it approximates.  Formally, if $r$ is an
exact real value and $[\ell, u]$ is its representation:
\[
  \ell \;\leq\; r \;\leq\; u
\]
Equivalently, rounding must always go \emph{outward}: the lower bound rounds
toward $-\infty$ (never up), the upper bound rounds toward $+\infty$ (never
down).  Rounding toward the inside of the interval is not permitted, as it
would exclude the true value and invalidate the IA-DEVS enclosure guarantee.

\subsubsection*{Consequences by TIME type}

\paragraph{\texttt{decimal<N>}.}
Values expressible as multiples of $10^{-N}$ are represented exactly; no
rounding occurs and intervals may be point intervals.  All experiments in this
repository that use \texttt{decimal<3>} with periods that are exact
multiples of $1\,\text{ms}$ satisfy this condition automatically.

\paragraph{Floating-point (\texttt{float}, \texttt{double}).}
IEEE~754 arithmetic does not guarantee outward rounding by default.  When
constructing an interval endpoint from a computed value, the caller must
use directed rounding:
\begin{itemize}
  \item Lower bound: round toward $-\infty$ (\texttt{fesetround(FE\_DOWNWARD)})
  \item Upper bound: round toward $+\infty$ (\texttt{fesetround(FE\_UPWARD)})
\end{itemize}
Experiments that use floating-point TIME must document whether directed
rounding is applied, and if not, what approximation error is being accepted.

\textbf{Compiler note.}  Many compilers apply floating-point optimizations
(reassociation, contraction, flush-to-zero) that silently violate the active
rounding mode set by \texttt{fesetround}.  To ensure directed rounding is
actually respected, disable these optimizations explicitly.  With GCC and
Clang, the relevant flags are \texttt{-frounding-math} (honours the runtime
rounding mode) and \texttt{-fno-unsafe-math-optimizations} (or equivalently,
avoiding \texttt{-ffast-math}).  Without these flags, interval endpoints may be
rounded in the wrong direction even when the correct rounding mode is set at
runtime.

\paragraph{Integer types.}
Integers are exact; intervals may always be point intervals when the modelled
quantity is an integer.

% ─────────────────────────────────────────────────────────────────────────────
\subsection*{Z-Function Coupling}

UA-DEVS and IA-DEVS, as defined in~\cite{VDW21}, couple components through
\emph{translation functions} $Z_{i,d}$, one per influencer pair $(i, d)$.
$Z_{i,d}$ maps the output of component $i$ to the input of component $d$;
the influencer set $I_d$ lists every $i$ that can send to $d$, and $I_N$
captures which components contribute to the coupled model's own output.
The $Z$-function mechanism is inherited from Zeigler's original DEVS coupled
formalism~\cite{Zeigler00}, carried into the uncertainty extensions of
UA-DEVS and IA-DEVS.

Ports are a feature of PDEVS-with-ports (the formalism Cadmium implements),
not of DEVS or its uncertainty extensions.  CadmIA implements the formalism
faithfully: each component pair $(i, d) \in I$ is registered with a
\texttt{translation\_fn} (the Z function), and the coordinator applies it
whenever $i$ fires and $d \in I_d$.

When all components share the same value type, $Z_{i,d}$ is often the
identity function; signal semantics are then distinguished by value rather
than by channel.  Each experiment spec must document its value-encoding
convention and the Z functions it uses.

% ─────────────────────────────────────────────────────────────────────────────
\subsection*{From Formalism to C++ Implementation}

The IA-DEVS atomic functions --- $\mathrm{TA}$, $\Lambda$,
$\Delta_{\text{int}}$, $\Delta_{\text{ext}}$ --- are defined over interval
states and interval values, but they are \emph{not called directly} by the
simulator engine.  The simulation algorithms from~\cite{VDW21} (the Bound
procedures) interleave these functions with interval bookkeeping: computing
$t_{\text{next}}$ as $t_{\text{last}} \oplus \mathrm{TA}(s)$, resolving
simultaneous imminence, and routing outputs through $Z$ functions before
calling $\Delta_{\text{ext}}$.

What the C++ atomic model actually provides is the \emph{raw} $\mathrm{TA}$,
$\Lambda$, $\Delta_{\text{int}}$, $\Delta_{\text{ext}}$ as member functions.
The simulator (CadmIA's \texttt{simulator} and \texttt{coordinator}) combines
them in one Bound step during \texttt{engine\_star} / \texttt{engine\_x}.
Model authors therefore implement the four atomic functions; the combined
Bound logic is the engine's responsibility.

% ─────────────────────────────────────────────────────────────────────────────
\subsection*{SELECT Function}

When multiple components are simultaneously imminent, the CadmIA coordinator
calls a user-supplied \textsc{select} function to pick the next component to
advance.  The choice affects the simulation trace whenever simultaneous events
carry causal dependencies.  Each experiment spec must state the \textsc{select}
policy and justify it.
